\section{Scope, Edition, and Qualifications}

\subsection*{The object}

This study treats the \emph{Order of Mass} --- the \latin{Ordo Missae}, nn.\ 1--146 --- of the
Roman Rite, together with the General Instruction that governs it and the Appendix to the Order
of Mass so far as the 2008 emendations touch it. ``Order of Mass'' here means the whole ordered
sequence of the Mass with a congregation, not the sung Ordinary (\latin{Kyrie, Gloria, Credo,
Sanctus, Agnus Dei}) and not the complete Missal. Stable does not mean invariant: the study
distinguishes throughout between elements that occur at every celebration, elements conditioned
by day or rank, elements the Missal offers as alternatives, and elements determined by
territorial adaptation.

Excluded: the temporal and sanctoral Propers; the Lectionary and its text; the Eucharistic
Prayers for Reconciliation and for Various Needs beyond the 2008 corrections that touch them;
the Eucharistic Prayers for Masses with Children beyond the fact of their removal from the 2008
Latin Missal; concelebration, Masses without a congregation, Masses with the Bishop, ritual and
votive Masses; the Order of Mass of the 1962 Missal except as a comparand at named points; and
every question of the present canonical discipline governing the older Missal.

\subsection*{The editions}

\begin{itemize}[itemsep=0.2em]
\item \textbf{Controlling Latin rite:} \latin{Missale Romanum}, \latin{editio typica tertia},
\latin{reimpressio emendata} (Vatican City, 2008). No exact artifact of the 2008 altar book was
obtained or is registered. The study therefore works from two controls in combination: the
registered secondary digital reproduction of the 2002 third typical edition, and the official
published list of the 2008 variations. Every statement about the 2008 state is either supported
by the variation list or is a statement that the variation list records no change at that place;
the latter is bounded negative evidence, since the 2008 notice states that further corrections of
accent, punctuation, colour, and typography were made without being listed.
\item \textbf{Latin General Instruction:} the \latin{Institutio Generalis Missalis Romani} as
printed in the same 2002 reproduction, with the one 2008 emendation at n.\ 149 noted in place.
\item \textbf{English:} \emph{The Roman Missal, Third Edition, for use in the Dioceses of the
United States of America} (2011), and the General Instruction in the United States English state
emended through 2021. Chapter II (nn.\ 27--90) of that state was obtained and read. Chapters III
and IV were not obtainable: the host returned an interstitial challenge page rather than the
chapter. Nothing in this study rests on nn.\ 160 or 281--287, and the study says so where the
gap affects a movement.
\item \textbf{Comparand:} \latin{Missale Romanum}, \latin{editio typica} (Vatican City, 1962),
in the registered facsimile, used at the Canon and at named points of comparison only.
\end{itemize}

\subsection*{Territory and date}

All statements about territorial adaptation are for the dioceses of the United States of
America and are stated as of 25 July 2026. The 2011 English is the state in use since the First
Sunday of Advent 2011; the General Instruction state is that carrying the 2021 emendations. This
study did not inspect the confirmation decrees themselves and cites their dates from the
registered official record of the United States General Instruction.

\subsection*{Rights and the treatment of texts}

The Latin \latin{editio typica} is quoted where the argument requires it, and the Latin
quotations in this study are evidence, not texts offered for recitation. The approved English
translation of the Order of Mass --- the people's parts, the orations, and the Eucharistic
Prayers --- is under copyright and is \emph{not} reproduced here in any form. Where the argument
turns on a short English phrase whose exact wording is the point at issue (``and with your
spirit,'' ``consubstantial with the Father,'' ``for many,'' ``under my roof''), that phrase is
quoted and no more. Nowhere is the approved English paraphrased closely enough to reconstruct
it.

That constraint has consequences, and they are stated where they bite rather than in a blanket
disclaimer: the Collect and the Prayer over the Offerings are analysed structurally without
English, and this is said at each; the Prayer after Communion likewise; the four dismissal
formulas are given in Latin only; the 2011 English of the Roman Canon and of the Eucharistic
Prayers is not assessed at all, and section 11 explains what that leaves undone. The Lectionary
text is under copyright and is never reproduced. Public-domain English for Scripture, where
needed, would be the Douay--Rheims version held in this repository; in the event, no scriptural
passage required quotation in English, and the biblical loci are cited rather than quoted.

Patristic and historical Latin quoted in this study --- Ambrose's \latin{De sacramentis} from
\emph{Patrologia Latina} XVI (1880), and the Verona palimpsest text from Hauler's Teubner
edition of 1900 --- is public domain in the United States by age of publication; the page images
consulted are faithful photographic reproductions adding no original authorship.

\subsection*{Evidence classes and their limits}

Claims in this study fall into four classes and are distinguishable in the notes.

\emph{Verified source text} --- wording read in an identified witness. This covers the Latin
Order of Mass and General Instruction (read in the registered 2002 reproduction), the 2008
variation list and notice (read at the official Dicastery PDF, hash-matched to the registered
artifact), the 2006 circular letter on \latin{pro multis} (read at the official Dicastery PDF),
\emph{Liturgiam authenticam} (read at the Holy See portal), the United States General
Instruction Chapter II (read at the official USCCB presentation), Paul VI's constitution
\latin{Missale Romanum} in Latin, the 1962 Canon (read at rendered page images of the registered
facsimile), Ambrose \latin{De sacramentis} IV (read at page images), and the Verona anaphora
(read at page images).

\emph{Checked citation} --- a claim checked at an exact locus, such as the dates of the 2010
confirmations, taken from the registered official record rather than from the decrees.

\emph{Source-grounded synthesis} --- the comparative tables, the structural readings of the
movements, and the whole-action account. These are editorial conclusions drawn from the checked
evidence and are labelled as such where they could be mistaken for the sources' own statements.

\emph{Editorial proposal} --- the Project synthesis block and its counterargument. These are
this study's own judgements, marked in place.

Specific limitations that a reader should carry away:

\begin{itemize}[itemsep=0.15em]
\item The controlling Latin artifact is a secondary digitally typeset reproduction, not a page
facsimile of the printed altar book. Its typesetting carries visible defects. Wording claims are
made only where the reading is unambiguous, and anyone needing an exact reading should consult a
printed altar book.
\item No exact 2008 altar-book artifact was obtained. Every 2008 claim is mediated by the
published variation list.
\item General Instruction chapters III--VIII were not read in full; the study's negative claims
about what the Missal does and does not require are bounded accordingly, and this bound is
restated inside the Project synthesis where it matters.
\item Prayer IV's relation to the Antiochene anaphoral family is asserted at the level of form,
from the Missal's own text and rubric. No edition of the anaphora of St Basil was examined and
no verbal dependence is claimed.
\item The authorship, date, place, and redactional unity of the church order conventionally
called the \latin{Traditio apostolica} are disputed. This study uses the Verona Latin as a
witness to a text and makes no attribution to Hippolytus.
\item The drafting history of the new Eucharistic Prayers is not established here. Promulgation
by papal act is documented; individual authorship is not claimed.
\item Reception history, celebration practice, music, architecture, and the sociology of the
reform are outside the study's scope entirely.
\end{itemize}

\subsection*{Review state}

This study has received internal source and production review only. It has not received
independent review by a liturgical historian, a patristics specialist, a Latinist, or a
sacramental theologian; it carries no imprimatur, nihil obstat, or ecclesiastical approval; and
it is not an altar book, a ritual text, a critical edition, or a substitute for the liturgical
books themselves.
